\documentclass[12pt,a4paper]{article} 
\usepackage[margin=1in]{geometry} 
\usepackage[T1]{fontenc} 
\usepackage{lmodern} 
\title{Informácie o projekte}
\author{Diana Horpinchenko, Anhelika Kcharchenko, Lev Hrabovyi, Yehor Hrebnov} 
\begin{document} 
\maketitle \section*{Názov projektu} 
\textbf{Názov:} Memory Integration \& Mind Imitation Creator
 \vspace{0.5cm}
 \textbf{Acronym projektu:} MIMIC 
\vspace{0.5cm} 
\section*{Popis projektu} 
Mimic je aplikácia, ktorá dokáže vytvoriť digitálnu verziu človeka na základe jeho spôsobu komunikácie. Naša aplikácia pracuje s textami, ktoré používateľ nahrá do systému, napríklad so starými správami, chatmi alebo e-mailami. Z týchto údajov sa postupne vytvorí profil konkrétnej osoby, ktorý nazývame „Mimic“.

Na začiatku používateľ jednoducho vyberie, koho chce vytvoriť, a nahrá dostupnú históriu komunikácie. Systém potom analyzuje, ako daný človek píše. Sleduje napríklad používané slová, dĺžku viet, slang, emotikony, spôsob vyjadrovania alebo to, ako reaguje na rôzne témy. Čím viac komunikácie aplikácia má k dispozícii, tým presnejšie dokáže napodobniť charakteristický štýl danej osoby.

Po vytvorení Mimica s ním môžeme normálne komunikovať cez chat. Aplikácia generuje odpovede podľa vytvoreného profilu, takže konverzácia pôsobí podobne ako komunikácia s pôvodnou osobou. Pred samotnou konverzáciou môžeme nastaviť aj náladu alebo emóciu Mimica. Môžeme napríklad skúsiť, ako by komunikácia vyzerala, keby bol človek nahnevaný, smutný alebo naopak dobre naladený.

Jednou z možností využitia je príprava na náročné rozhovory. Ak napríklad chceme požiadať šéfa o zvýšenie platu, môžeme si najprv vytvoriť jeho Mimica a skúsiť si s ním rozhovor. Aplikácia nám tak umožní vyskúšať rôzne situácie a reakcie ešte pred skutočným rozhovorom.

Mimicry môže fungovať aj trochu iným spôsobom. Používateľ môže nahrať napríklad svoje vlastné staré správy a vytvoriť si Mimic svojho minulého ja. Potom sa s ním môže rozprávať a porovnať, ako sa jeho spôsob komunikácie a názory časom zmenili.

Dôležitou súčasťou aplikácie je aj ochrana údajov. Keďže pracujeme so súkromnými správami, dáta musia byť šifrované a spracované iba na účely fungovania aplikácie. Používateľ má zároveň kontrolu nad tým, aké údaje do systému nahrá a aký profil vytvorí.

Mimicry teda funguje ako nástroj na simuláciu komunikácie. Namiesto obyčajného chatbota vytvára profil konkrétneho spôsobu komunikácie a umožňuje nám skúšať rôzne rozhovory v interaktívnom prostredí.
\end{document}